% Table 4.3 (dlsu-manuscript/chapter_4.tex, tab:interpretation), wording unchanged.
\documentclass[12pt,border=10pt]{standalone}
% Shared style for the slide figures: sans-serif text (Helvetica, close to the Arial of the charts),
% the same colors as charts.py, and the manuscript's table tools.
\usepackage[T1]{fontenc}
\usepackage[scaled]{helvet}
\renewcommand{\familydefault}{\sfdefault}
\usepackage{amsmath,amssymb}
\usepackage{sansmath}
\usepackage{xcolor}
\usepackage{booktabs,array}
\newcolumntype{L}[1]{>{\raggedright\arraybackslash}p{#1}}
\usepackage{tikz}
\usetikzlibrary{arrows.meta,positioning,calc}

\definecolor{ink}{HTML}{0B0B0B}
\definecolor{inktwo}{HTML}{52514E}
\definecolor{muted}{HTML}{898781}
\definecolor{gridc}{HTML}{E1E0D9}
\definecolor{basec}{HTML}{C3C2B7}
% one color per selection rule, in the fixed categorical order of the charts
\definecolor{cp}{HTML}{2A78D6}     % k-Changepoint (content)
\definecolor{sme}{HTML}{EB6834}    % SME (position)
\definecolor{posc}{HTML}{1BAF7A}    % positional (position)
\definecolor{rnd}{HTML}{EDA100}    % random (chance)
\definecolor{bocpd}{HTML}{E87BA4}  % k-BOCPD (content), categorical slot 5
\definecolor{temporal}{HTML}{6250D6}  % temporal information added to the mean (frame combinations figure)

\newcommand{\chip}[1]{\tikz[baseline=0.1ex]\fill[#1, rounded corners=1pt] (0,0) rectangle (0.32cm,0.32cm);}
\AtBeginDocument{\sansmath\color{ink}}
\hyphenpenalty=10000 \exhyphenpenalty=10000   % no hyphenated words on slides

\begin{document}
\renewcommand{\arraystretch}{1.35}
\begin{tabular}{@{}l@{}}
   \begin{tabular}{@{}L{5.6cm}L{10.4cm}@{}}
      \toprule
      Outcome at a frame budget $k$ & Interpretation \\
      \midrule
      \multicolumn{2}{@{}l}{\textit{Each content-based rule against positional and random selection}} \\
      Content $>$ random and content $>$ positional & The frames this rule selects carry information beyond both an arbitrary and an evenly spaced choice of frames. \\
      Content $\approx$ random, both $>$ positional & Evenly spaced frames are a poor choice, but this rule is not shown to add information beyond an arbitrary choice. \\
      Positional $>$ content & Covering the whole clip evenly is more useful than the frames this rule selects. \\
      Content $\approx$ random $\approx$ positional & No evidence that the selection rule matters at this frame budget. \\
      \midrule
      \multicolumn{2}{@{}l}{\textit{The two content-based rules against each other}} \\
      \chip{bocpd}\ $k$-BOCPD $>$ \chip{cp}\ $k$-Changepoint & The beginnings of lasting changes are more informative frames than the moments of fastest movement. \\
      \chip{cp}\ $k$-Changepoint $>$ \chip{bocpd}\ $k$-BOCPD & The moments of fastest movement are more informative frames than the beginnings of lasting changes. \\
      \chip{cp}\ $k$-Changepoint $\approx$ \chip{bocpd}\ $k$-BOCPD & The two definitions of content select equally useful frames at this frame budget. \\
      \midrule
      Any other pattern & Reported as observed, without a general conclusion about the selection rules. \\
      \bottomrule
   \end{tabular}\\[4pt]
   \parbox{16cm}{\small\raggedright\textcolor{inktwo}{Frames combined by their mean. ``Content'': either content-based rule, each read separately. ``$>$'': a reliable difference (significant after correction; in all five seeds or draws where these apply). ``$\approx$'': equivalent within $\pm 2$ accuracy points.}}
\end{tabular}
\end{document}
